\section{问题重述}

\subsection{问题背景}

研究对象为由居民负载、光伏电源、储能设备和外部电网构成的非孤岛微网。小区用电必须持续得到满足，但光伏具有明显的日内波动，外网电价也随交易时段变化。储能可以把低价时段或光伏富余时段的电量保存到后续使用。一天内各时段由储能状态相互联结：当前的充放电动作会改变后续储电量，进而影响以后能够少买多少电、是否需要紧急购电以及全天总费用。

题目以 $10$ min 为基本调度粒度。储能额定容量为 $12000$ kWh，允许的储电量范围为 $1200\sim10800$ kWh，最大充、放电功率均为 $5000$ kW，充放电效率均按 $90\%$ 处理。于是，一天 $144$ 个时段通过储能状态递推形成连续的决策链：某一时段多充电，会提高后续可利用的储能；某一时段过早放电，则可能使真正高价或缺电时段失去调节余量。

除上述物理联系外，四个问题还存在一条更重要的递进主线，即\textbf{不同决策时刻能够获得的信息并不相同}。问题一给出的信息最完整，可以直接制定确定性全天计划；问题二要求计划先于当天实际源荷实现，并为预测偏差设置高价紧急购电；问题三允许在日内收到新预报后修改尚未执行的计划；问题四再把固定的典型日电价推广为逐日波动价格。由此，四问本质上是在同一微网物理系统上，逐步增加不确定性、信息更新和结算约束。

\subsection{问题提出}

\indent\textbf{问题一：}在典型日电价、负载和光伏预测均给定的条件下，于每日 $0{:}00$ 一次性确定全天各 $10$ min 时段的计划购电量和储能充放电量，使所有时段均满足负载，并保证日初、日末储电量相同。在论文中给出题目指定时段的购电结果、全天购电量与费用，以及各 $4$ h 时段的储能充放电量；完整方案写入 \texttt{result1.xlsx}。

\indent\textbf{问题二：}当全年负载和光伏随日期变化时，仍要求每天 $0{:}00$ 先制定全天计划；若实际运行后发现供电不足，缺口需按交易时刻正常电价的 $5$ 倍紧急购入。这里存在两个相互制约的问题：计划过少会放大高价紧急购电风险，计划过多又会增加正常购电费用。模型需利用历史数据形成可执行的日前计划，并在实际偏差出现后合理调用储能缓冲。按题目要求给出 2025-03-20、2025-06-21、2025-09-23、2025-12-21 四个指定日期的购电、储能和紧急购电结果，完整方案写入 \texttt{result2.xlsx}。

\indent\textbf{问题三：}每天 $0{:}00$、$6{:}00$、$12{:}00$、$18{:}00$ 均可获得未来 $24$ h 的光伏预报。$0{:}00$ 的全天计划是后续调整的基准，新的预报到达后只能修改尚未执行的部分，并对向上、向下调整分别计价。因而需要建立状态连续的滚动优化模型，并通过对照实验回答：新增的日内预报是否值得用于修改购电计划、各更新节点的经济价值有多大。同样对 2025-03-20、2025-06-21、2025-09-23、2025-12-21 给出题目要求的计划购电、调整购电、储能和紧急购电结果，完整结果写入 \texttt{result3.xlsx}。

\indent\textbf{问题四：}在问题二、三已有调度机制上进一步引入逐日波动电价。根据附件4重新决定购电和储能动作，比较固定典型日电价与实时波动价格下的调度差异，并对上述四个指定日期给出相应重算结果，分别形成与问题二、问题三对应的完整结果，写入 \texttt{result4-2.xlsx} 与 \texttt{result4-3.xlsx}。

\subsection{研究现状综述}

微网经济调度问题在数学上通常被建模为以购电费或综合运行成本最小为目标的线性规划或混合整数规划，外网购电、分布式电源消纳与储能充放电被统一纳入同一组能量平衡约束中 \cite{wu2014dynamic,nemati2018optimization}。确定性模型适用于负载、电价与光伏均已知的日前场景，其优点是约束结构清晰、求解稳定可靠。

当光伏出力存在预测误差时，主流处理方式有三类：一是\textbf{情景法／两阶段随机规划}，用有限个历史或生成情景刻画不确定性，第一阶段确定与情景无关的日前计划，第二阶段按情景确定执行量 \cite{wu2016solution,sun2025distributionally}；二是\textbf{鲁棒优化}，在最坏情景下保证可行性，代价是结果偏保守 \cite{zhang2018robust}；三是\textbf{滚动时域优化（模型预测控制）}，在新信息到达后仅重优化剩余时域并只执行近期决策，天然适配“预报分批发布”的决策结构 \cite{elkazaz2020energy,silva2020optimal}。含光伏与储能的日前调度研究普遍指出，储能的价值高度依赖预测精度与调度时域长度 \cite{luo2020optimal}，而储能的运行成本又与其循环深度密切相关 \cite{maheshwari2020optimizing}。

针对“预测值应偏向哪一侧”的问题，报童模型给出的结论是：当高估与低估的边际损失不对称时，最优点决策由相应分位数确定 \cite{qin2011newsvendor}。这一思想已被用于电力系统中的发电商投标策略与偏差电量定价 \cite{polat2024renewable}。在预测方法层面，基于数值天气预报的机器学习模型是目前光伏功率预测的重要方法 \cite{markovics2022comparison}；在市场机制层面，实时电价与需求响应的耦合会显著改变用户的购电行为 \cite{wang2015review}。

现有研究多聚焦于算法本身，对\textbf{信息集边界}与\textbf{结算规则}的刻画相对有限。在业务未提供全天计划负荷时，日前优化若直接使用当天真实负载，会混入理想信息优势；若实际业务已提供该信息，则可将其作为信息更充分的对照。比较不同策略时也应尽量保持预测方法、约束和费用规则中的非研究因素一致，才能解释费用变化来自何处。本文据此把信息边界作为建模原则，并用单因素对照与配对控制支持归因；求解则利用线性／混合整数线性结构，避免引入不必要的随机搜索误差。
